\begin{abstract}
A run-to-run standard error for a benchmark average built from each benchmark's variance alone ignores the covariance between benchmark scores. Our estimator SNAP splits each benchmark's items into two halves and takes cross-half products of run-centred scores so item-sampling noise drops out, and we state the conditions for identifying run covariance and calibrate our recipe-cluster intervals in simulation. As a worked example we apply SNAP to 375 released DataDecide runs spanning 125 configurations and ten benchmarks. On per-byte margins the battery average spreads more than independence predicts, with a standard deviation 1.244 times the independent value (interval \ci{1.143}{1.338}, lower endpoint 1.095 under a test inversion added later), which makes the ten-benchmark average as noisy as an average of 6.5 independent benchmarks with the same average variance, and every recipe, recipe-pair and benchmark deletion leaves an interval that excludes one. \outcome{\Ronecase}{1}{R1 pass}{\outcome{\Rtwocase}{1}{R2 not supported}{Under rules fixed before scoring, nine seeds at five PolyPythias sizes give margin inflation of \pending{R1 lambda} with interval \ci{\pending{lo}}{\pending{hi}}, and pairing that battery with a disjoint item bank leaves \pending{cross lambda} with the shared-item explanation rejected.}\outcome{\Rtwocase}{2}{R2 supported}{Under a rule fixed before scoring, nine seeds at five PolyPythias sizes give margin inflation of \pending{R1 lambda} with interval \ci{\pending{lo}}{\pending{hi}}, although a disjoint item bank supports a shared item component.}\outcome{\Rtwocase}{3}{R2 undecided}{Under a rule fixed before scoring, nine seeds at five PolyPythias sizes give margin inflation of \pending{R1 lambda} with interval \ci{\pending{lo}}{\pending{hi}}, although a disjoint item bank leaves the source of that inflation undecided.}}\outcome{\Ronecase}{2}{R1 fail}{On a second family, nine seeds at five PolyPythias sizes give margin inflation of \pending{R1 lambda} with interval \ci{\pending{lo}}{\pending{hi}}, which fails a rule fixed before scoring at a design with 40 within-size contrasts.} We find that the size of the factor belongs to the battery. BoolQ carries 68\% of the margin covariance trace, and removing it still leaves 1.786 across the nine remaining benchmarks. However, on accuracy our interval \ci{0.993}{1.157} includes one under a low-powered design. A registered test on four unseen tasks fails its lower-limit rule at 1.217 with interval \ci{0.872}{1.498}, while on the same runs the original ten benchmarks give 1.240 with \ci{1.054}{1.397}. Every other analysis is exploratory, no registration carries an independently corroborated timestamp, and on real recipe comparisons the covariance changes the wrong-call rate by at most 0.009. We recommend paired replicate scores across the chosen battery with a benchmark-removal check.
\end{abstract}

\section{Introduction}
A benchmark average can change across training runs even when the evaluation items stay fixed. For an equally weighted average of $K$ benchmarks, independent run deviations give variance $K^{-2}\operatorname{tr}\Sig$, where $\Sig$ is their cross-benchmark covariance matrix. Covariance adds the off-diagonal terms to that variance, and assuming independence therefore weakens a comparison's error control. In simulation with margin-like covariance, such a comparison declares a difference in 0.113 of replicates with a true gap of zero, against a nominal rate of 0.05.

We estimate item-general covariance from two disjoint item halves per benchmark, centring both half scores across the replicate runs within each configuration and taking their cross-half products. Under zero-mean item noise and conditional independence across halves, they estimate run covariance, including its diagonal. We adapt the established repeated-measurement construction to this setting and call it SNAP, Seed Noise Across Phenotypes, where phenotype names the score we call a trait. We keep the name though the released runs differ in checkpoint step, training budget and seed.

Our empirical study uses DataDecide \citep{magnusson2025}, with 25 data recipes and five model sizes that each have three released replicates, and those 125 configurations supply 375 runs across ten benchmarks. Our primary analysis reads released item scores and requires no model inference, and we compare per-byte log-likelihood margins with accuracy on the same items. These two scales respond differently to changes in confidence.

We find the scales give different inflation estimates. Margin inflation is 1.244 with recipe-cluster interval \ci{1.143}{1.338}, which excludes one, while accuracy gives 1.078 with interval \ci{0.993}{1.157}, which includes one and bounds accuracy inflation at 1.157 in the primary analysis (Table~\ref{tab:primary}).

The accuracy interval no longer includes one under five of the 25 recipe removals, removal of the 750M band, and removal of WinoGrande or BoolQ, while a simulation at the fitted accuracy cell excludes one in 0.111 of replicates and would have detected a true ratio of 1.10 in only 0.135 of them. We read the full-battery accuracy estimate as a bound and not as evidence of independence.

Because we built our estimator on these ten benchmarks, we registered a test on SciQ, MedMCQA, and multiple-choice DROP and CoQA, scored on the same checkpoints at 530M, 750M, and 1B, with a registered rule requiring the lower interval limit for held-out margin inflation to exceed one. However, the registered test fails on all 225 runs, where the held-out estimate is 1.217 with interval \ci{0.872}{1.498}. The original ten benchmarks on the same runs give 1.240 with interval \ci{1.054}{1.397} and pass it, a comparison fixed before any held-out score existed (Section~\ref{sec:heldout}).

A single model family and a seed label that also changes the training batch leave open whether the inflation belongs to DataDecide's design. We therefore scored all 45 PolyPythias runs at five sizes on the same 37,682 items, where each of nine seeds reruns one recipe, and paired that battery with a disjoint bank of 6,808 items built at the same OLMES commit. Two rules fixed before scoring read the results, and Section~\ref{sec:family} reports that margin inflation there is \pending{R1 lambda} with interval \ci{\pending{lo}}{\pending{hi}} and that the shared-item explanation is \outcome{\Rtwocase}{1}{R2 not supported}{rejected}\outcome{\Rtwocase}{2}{R2 supported}{supported}\outcome{\Rtwocase}{3}{R2 undecided}{undecided}.

This paper makes four contributions to measuring run covariance. We define the item-general run covariance that cross-half products estimate, and state which assumptions the released runs can't be shown to satisfy. Algorithm~\ref{alg:snap} gives the measurement recipe with its interval calibration, where wild recipe-cluster coverage runs from 0.928 to 0.954 across twelve simulated populations. We measure that covariance on two score scales across the 375 released runs, where margin inflation is 1.244 with an interval that excludes one under every recipe, pair and benchmark deletion, at a magnitude that BoolQ sets. At that value the ten benchmarks average like 6.5 independent ones of the same average variance, with an interval from 5.6 to 7.7. The accuracy interval includes one at a design that would detect a true ratio of 1.10 in only 0.135 of replicates.

We register a test on unseen tasks, and it fails at its registered scope, where the held-out margin interval runs 0.626 wide against 0.343 for the original ten benchmarks on the same runs and the rule reads the lower limit. We show what the covariance changes in practice, where a correction estimated on the same battery returns the simulated false-positive rate from 0.113 to 0.051 and leaves 0.101 at a battery whose own ratio is 1.5. On real recipe comparisons at four sizes it changes the wrong-call rate by at most 0.009 with every interval for that change reaching zero. Neither the analysis plan nor the protocol's registration carries an independently corroborated timestamp, and the plan's screening split was never formed. Every analysis apart from the registered test is exploratory.

